\newpage
\section{Introduction}
\label{ch:introduction}

\subsection{Background of the Study}
\label{sec:background}

Assessing rice pollen viability is important for understanding the reproductive potential of rice plants and supporting research in plant breeding. However, conventional pollen viability assessment commonly involves staining pollen grains and manually examining and counting them under a microscope, a process that can be labor-intensive and time-consuming, particularly when large numbers of grains must be evaluated~\cite{tan2022,ascari2020}. To address the need for automated assessment, Go et al.~\cite{go2019} developed an image-analysis protocol for identifying viable and nonviable pollen in selected rice genotypes. Their method employed edge detection, dilation, region filling, and $k$-means clustering to identify individual pollen grains and classify their viability based on grayscale intensity. Although the method demonstrated promising results, its performance was affected by the presence of overlapping or clustered pollen grains. As reported in the study, accuracy decreased from 98.3\% for images without lumped pollen to 85.06\% for images containing lumped pollen~\cite{go2019}. This limitation suggests that conventional image-processing techniques may struggle to distinguish individual grains when their boundaries are difficult to identify.

Subsequent studies have explored machine learning and deep learning to improve automated pollen detection, segmentation, and viability classification. While standard object detection models can locate pollen grains using bounding boxes~\cite{tan2022}, severely overlapping or clustered grains require instance segmentation to generate precise pixel-level masks for each individual grain~\cite{zhang2023}. These learning-based approaches demonstrate strong potential for automating pollen analysis under challenging image conditions. Nevertheless, based on the literature reviewed in this study, the effectiveness of these methods for classifying rice pollen viability under overlapping or clustered conditions remains insufficiently established. In particular, it remains necessary to determine whether a model can distinguish individual rice pollen grains and classify their viability reliably when their boundaries overlap, and whether its performance under these conditions improves upon that of the existing image-analysis protocol. Addressing this gap requires an approach that can learn relevant visual features directly from annotated images rather than relying solely on manually selected image-processing parameters.

To address this problem, this study proposes fine-tuning a lightweight, pretrained YOLOv8 nano instance-segmentation model (YOLOv8n-seg) to detect individual rice pollen grains, generate a separate segmentation mask for each detected grain, and classify each grain as viable or nonviable. The nano variant provides a low computational footprint suitable for resource-constrained environments while mitigating the risk of overfitting on a small dataset. The proposed model will be evaluated against the image-analysis protocol of Go et al.~\cite{go2019} using stained microscopic images from six rice genotypes. Both methods will be assessed separately on images with and without lumped pollen to determine how overlapping grains affect their performance. Detection and classification metrics, together with comparisons between automated viability estimates and manual counts, will be used to evaluate the proposed approach. Through this comparison, the study seeks to determine whether instance segmentation offers a more reliable approach to automated rice pollen viability classification, particularly in images containing overlapping or clustered grains. The specific scope, limitations, and practical significance of this investigation are detailed in the subsequent subsections of this chapter.

\subsection{Scope and Limitations}
\label{sec:scope-limitations}

This study focuses on the automated classification of rice pollen viability from stained microscopic images. It will use images from six rice genotypes: IR68652-3B-8-1, Sadri, CSR28, IR67075-2B-5-2, Rc222, and Pokkali, subject to permission to access the dataset used by Go et al.~\cite{go2019}. If additional images are required, samples will be collected from the same genotypes using comparable staining and imaging procedures. The study will fine-tune a pretrained YOLOv8n-seg instance-segmentation model to identify individual pollen grains and classify them as viable or nonviable. Its performance will be evaluated on images with and without lumped pollen and compared with the image-analysis protocol of Go et al.~\cite{go2019}, which uses edge detection, dilation, region filling, and $k$-means clustering.

The evaluation will focus on the model's ability to detect and classify individual grains, the effect of overlapping or clustered pollen on performance, and the agreement between automated viability estimates and manual counts. Stratified five-fold cross-validation will be used to evaluate the methods, with precision, recall, and F1-score used for detection and classification, and mean absolute error (MAE) and Pearson correlation used to assess the agreement of estimated viability percentages with manual counts. The comparison will be limited to the proposed YOLOv8n-seg model and the reimplemented protocol of Go et al.; other deep-learning architectures will not be evaluated in this study.

Several limitations may affect the findings. First, the study depends on access to the original dataset, and the number and diversity of available images may limit the range of conditions represented. The images cover six rice genotypes and may not represent all rice varieties, staining procedures, microscopes, or image-acquisition conditions. Second, the model's performance will depend on the quality and consistency of grain annotations, especially when individual grains overlap and their boundaries are difficult to identify. Third, the baseline comparison may be affected by implementation choices where details of the original protocol are unavailable; the assignment of its $k$-means clusters to viable and nonviable classes will therefore be checked against the annotated images. Finally, the study evaluates image-based, stain-based classification and estimated viability percentages. It will not directly measure pollen germination, fertilization success, crop yield, or the model's performance in field conditions. These limitations will be considered when interpreting the results and determining the extent to which the proposed approach can be generalized.

\subsection{Significance of the Study}
\label{sec:significance}

This study may contribute to the development of more reliable automated methods for assessing rice pollen viability, particularly when pollen grains overlap or occur in clusters. Previous research has demonstrated the potential of image analysis and deep learning for pollen identification and viability assessment, but performance can vary with pollen type and image conditions~\cite{go2019,zhang2023,tan2022,li2026pvdt}. By evaluating a YOLOv8n-seg model on stained rice pollen images, this study will provide evidence on the usefulness and limitations of instance segmentation for this specific task.

For rice researchers and plant breeders, an effective automated method could help reduce the amount of manual counting needed when examining large numbers of pollen grains. By producing grain-level classifications and estimating the proportion of viable pollen, the proposed approach may offer a useful supporting tool for pollen assessment. It is intended to complement, rather than replace, established laboratory procedures; its practical value will depend on the results of the evaluation and the quality of the available images.

For researchers in computer vision and agricultural applications of artificial intelligence, this study will examine how a pretrained instance-segmentation model performs on a specialized microscopy task with overlapping objects. The comparison with the image-analysis protocol of Go et al.~\cite{go2019}, together with separate evaluation of images with and without lumped pollen, will help clarify whether the proposed approach addresses the limitations identified in the existing method. The use of cross-validation and comparison with manual counts will also provide a structured basis for assessing detection, classification, and viability-estimation performance.

Finally, the findings may serve as a reference for future studies on automated pollen analysis. If the proposed method proves effective, subsequent work could investigate additional rice genotypes, larger and more varied image datasets, other segmentation models, or validation against complementary biological measures. If its performance remains limited, the results may still identify challenges that need to be addressed before deep-learning-based pollen viability assessment can be applied more broadly.
